\chapter{МГД-стійкість і операційні межі}\label{ch:stability}

Рівновага з розд.~\ref{ch:equilibrium} корисна лише тоді, коли вона стійка. Сформулюємо енергетичний принцип, розберемо вертикальну нестійкість і пилку, далі~--- межі за $\beta$ і густиною, ELM, тиринг- і альфвенівські моди.

\section{Лінійна ідеальна МГД і енергетичний принцип}\label{sec:07-energy}

\begin{outofcorpus}[силовий оператор, $\delta W$ і джерела нестійкості]
Зсунемо кожен елемент плазми з рівноваги \eqref{eq:03-force} на мале зміщення $\vect{\xi}(\vect{r},t)$ (жирне $\vect{\xi}$; не плутати з гвинтовим кутом $\xi$ розд.~\ref{ch:particles}). Лінеаризовані ідеальні рівняння дають
\begin{equation}\label{eq:07-force}
\begin{gathered}
\rho_m\frac{\pa^2\vect{\xi}}{\pa t^2}=\vect{F}(\vect{\xi})\equiv\frac{1}{\mu_0}\big[(\grad\times\vect{B}_1)\times\Bvec+(\grad\times\Bvec)\times\vect{B}_1\big]-\grad p_1,\\
\vect{B}_1=\grad\times(\vect{\xi}\times\Bvec),\qquad p_1=-\vect{\xi}\cdot\grad p-\gamma_{\mathrm a}p\,\grad\cdot\vect{\xi},
\end{gathered}
\end{equation}
$\gamma_{\mathrm a}=5/3$. Оператор $\vect{F}$ самоспряжений. Тому для мод $\vect{\xi}\propto\ee^{-\mathrm{i}\omega t}$ величина $\omega^2$ дійсна, і нестійкість означає $\omega^2<0$ з інкрементом $\gamma=\sqrt{-\omega^2}$. Звідси \emph{енергетичний принцип}: рівновага стійка тоді й лише тоді, коли $\delta W=-\tfrac12\int\vect{\xi}^*\cdot\vect{F}(\vect{\xi})\,\dd V\ge0$ ($\delta W=0$~--- маргінальна) для всіх допустимих $\vect{\xi}$. Шукати власні частоти для цього не потрібно, досить знайти пробну функцію з $\delta W<0$. В «інтуїтивній» формі для плазми
\begin{equation}\label{eq:07-dW}
\delta W_{\mathrm p}=\frac12\!\int\!\Big[\frac{|\vect{B}_{1\perp}|^2}{\mu_0}+\frac{B^2}{\mu_0}\big|\grad\cdot\vect{\xi}_\perp+2\vect{\xi}_\perp\cdot\vect{\kappa}\big|^2+\gamma_{\mathrm a}p|\grad\cdot\vect{\xi}|^2
-2(\vect{\xi}_\perp\cdot\grad p)(\vect{\kappa}\cdot\vect{\xi}_\perp^*)-J_\parallel(\vect{\xi}_\perp^*\times\unitb)\cdot\vect{B}_{1\perp}\Big]\dd V,
\end{equation}
$\vect{\kappa}=\unitb\cdot\grad\unitb$~--- кривина силової лінії. Перші три доданки невід'ємні: це згинання силових ліній (шир-альфвенівська хвиля), стиснення поля і звук. Дестабілізують лише два останні. \textbf{Тиск і кривина} ($\grad p\cdot\vect{\kappa}>0$, «погана» кривина на зовнішньому обводі) відповідають за \emph{перестановні} (interchange) і \emph{балонні} (ballooning) моди. Балонна мода витягнута вздовж $\Bvec$ ($\kpar\approx0$, мале згинання) і локалізована там, де кривина погана. Її параметри~--- шир $s=(r/q)\,\dd q/\dd r$ і нормований градієнт $\Pb=-2\mu_0Rq^2p'/B^2$ (у літературі~--- $\alpha$); перша межа стійкості при $s\gtrsim1$ приблизно $\Pb\approx0{,}6\,s$. Перестановну стійкість у циліндрі задає критерій Сюйдама $rB_z^2(q'/q)^2/(8\mu_0)+p'>0$, а в токамаку при $q>1$ середня кривина сприятлива (Мерсьє). \textbf{Паралельний струм} $J_\parallel$ дає \emph{кінк}-моди (kink), що гвинтово згинають шнур. Вакуумну область і стінку враховує окремий доданок $\delta W_{\mathrm v}\ge0$.
\end{outofcorpus}

\textbf{NOVA-W}~\cite[с.~7--9]{Ward1992} розв'язує саме задачу на власні значення для \eqref{eq:07-force}, але з вакуумом, що містить резистивні провідники та котушки зворотного зв'язку. Збурене поле у вакуумі аксіально-симетричне й задається збуреним полоїдальним потоком. У~\cite{Ward1992} для нього взято окрему літеру і рахують його повним, тобто там стоїть $2\pi\,\delta\psi$ з точністю до знака конвенції (там $\Bvec=\grad\tor\times\grad\psi+g\grad\tor$). Потік $\delta\psi$ у вакуумі складається з внесків струмів котушок і пасивних провідників (через функцію Гріна \eqref{eq:03-green}) та інтегралів Гріна по межах вакуумної області~\cite[с.~8, (7)]{Ward1992}. Тонку стінку товщиною $\delta_w$ з питомим опором $\eta$ описує стрибок нормальної похідної~\cite[с.~9, (9)]{Ward1992}:
\begin{equation}\label{eq:07-wall}
\big[\!\big[\hat{\vect n}\cdot\grad\delta\psi\big]\!\big]=\frac{\mu_0\delta_w}{\eta}\frac{\pa\,\delta\psi}{\pa t}.
\end{equation}
Це закон Ома $K_\tor=(\delta_w/\eta)E_\tor$ для поверхневого струму разом з $E_\tor=-R^{-1}\pa_t\delta\psi$ (знак залежить від орієнтації нормалі). Ідеальна стінка ($\eta\to0$) «заморожує» потік, а за $\eta\to\infty$ стінки немає. Межа плазми зв'язана з вакуумом через $\xi_\psi\equiv\vect{\xi}\cdot\grad\psi=\sum_m\xi_{\psi m}\ee^{\mathrm{i}m\Theta}=\mp\delta\psi$ ($\Theta$~--- полоїдальний кут Ward на межі плазми)~\cite[с.~9, (11)]{Ward1992}. Відхилення від жорсткого зсуву вимірюють відношеннями $\xi_{\psi2}/\xi_{\psi1}$, $\xi_{\psi3}/\xi_{\psi1}$, $\xi_{\psi4}/\xi_{\psi1}$. Струм котушки зворотного зв'язку пропорційний різниці потоків на двох петлях, симетричних відносно екватора, і швидкості її зміни~\cite[с.~8, (8)]{Ward1992}. У джерелі коефіцієнти $\alpha_m$, $\beta_m$ стоять при повному потоці; ми переписуємо через $\delta\psi$: $G_{\mathrm p}=2\pi\alpha_m$, $G_{\mathrm d}=2\pi\beta_m$ (відношення $G_{\mathrm d}/G_{\mathrm p}$ те саме), а петлі позначаємо $(R_o,Z_o)$ (у джерелі $X_o$):
\begin{equation}\label{eq:07-PD}
I_m=G_{\mathrm p}\big[\delta\psi(R_{o1},Z_{o1})-\delta\psi(R_{o2},Z_{o2})\big]+G_{\mathrm d}\,\frac{\dd}{\dd t}\big[\delta\psi(R_{o1},Z_{o1})-\delta\psi(R_{o2},Z_{o2})\big].
\end{equation} Нижче $g_{\mathrm p}$~--- нормований $G_{\mathrm p}$ (у джерелі $\alpha_g$; $\mu_0/2$, помножене на коефіцієнт в А/(Вб$\cdot$рад); нормування джерела, розмірність не пояснено)~\cite[с.~22]{Ward1992}.

\section{Вертикальна нестійкість і керування положенням}\label{sec:07-vde}

\begin{outofcorpus}[жорстка модель і роль стінки]
Для вертикального поля з індексом спаду $n_{\mathrm d}$ (п.~\ref{sec:04-shape}) з $\grad\times\Bvec=0$ маємо $\pa B_R/\pa Z=\pa\BZ/\pa R$, тож $B_R\approx-n_{\mathrm d}\BZ Z/R$. На жорстко зсунутий шнур діє сила $F_Z=-2\pi|\Ip\BZ|\,n_{\mathrm d}Z$. Для витягнутої плазми $n_{\mathrm d}<0$, і сила штовхає шнур далі від екватора. Без стінки інкремент альфвенівський ($\gamma^{-1}\sim$ мкс). Вихрові струми в провідній стінці зі сталою часу $\tau_w$ сповільнюють рух: у схемі «жорсткий шнур + одне коло стінки» $\gamma\tau_w\approx|n_{\mathrm d}|/(n_{\mathrm w}-|n_{\mathrm d}|)$, де $n_{\mathrm w}>0$~--- індекс стабілізації стінкою, що зростає з наближенням стінки. При $|n_{\mathrm d}|\to n_{\mathrm w}$ мода знову стає ідеальною. Отже, пасивна структура лише переводить нестійкість на повільний масштаб $\tau_w$, а стабілізує її активний зворотний зв'язок, швидший за $\tau_w$.
\end{outofcorpus}

\begin{corpusexample}[деформованість проти зворотного зв'язку: NOVA-W~\cite{Ward1992}]
\textbf{Пасивні пластини} (оптимальні кути~--- п.~\ref{sec:04-shape}). Власна функція підлаштовується під розміщення провідників: гармоніки $m=2,3$ змінюються так, щоб зменшити стабілізуючі вихрові струми~\cite[с.~16]{Ward1992}. Коли пластини стоять з внутрішнього боку, зміщення біля них мале, і плазма ніби «прослизає» (slip) довкола провідників~\cite[с.~16, 28]{Ward1992}.
\textbf{Активний зв'язок, PBX-M} (модифікований розряд \#226879: $\Ip=\SI{567.4}{кА}$, $R_0=\SI{1.635}{м}$, $a=\SI{0.308}{м}$, $\kappa_{95}=1{,}951$, $\Btor(0)=\SI{1.20}{Тл}$, $\qnf=2{,}51$)~\cite[с.~20]{Ward1992}. \emph{Внутрішні петлі} $(R_o,Z_o)=(1{,}255;\,\pm0{,}10)$~м: з ростом $g_{\mathrm p}$ інкремент спершу швидко падає. Далі, при $g_{\mathrm p}\gtrsim1{,}5$, гармоніка $m=2$ змінює знак, $\xi_{\psi4}/\xi_{\psi1}$ зростає приблизно вчетверо, а $\gamma$ застигає на маргінальній межі навіть при $g_{\mathrm p}=6$. Диференційна складова не допомагає, бо коливань немає ($\omega_r=0$)~\cite[с.~22, 24]{Ward1992}. Причина в тому, що нульовий контур збуреного потоку зсувається на петлі, і сигнал зникає~\cite[с.~23--24]{Ward1992}. \emph{Центральні зовнішні} $(1{,}64;\,\pm0{,}56)$~м (як в експерименті): без диференційної складової $\gamma$ падає лише до $\approx\SI{20}{с^{-1}}$ при $g_{\mathrm p}=1{,}5$, далі зростає частота коливань. Стабілізують $G_{\mathrm d}/G_{\mathrm p}=0{,}05$ і $0{,}10$~\cite[с.~24--25]{Ward1992}. \emph{Далекі зовнішні}: навіть при $G_{\mathrm d}/G_{\mathrm p}=0{,}10$ інкремент не опускається нижче $\approx\SI{4}{с^{-1}}$ і виходить на плато при $g_{\mathrm p}\approx1{,}75$~\cite[с.~27]{Ward1992}. Для відношення $G_{\mathrm d}/G_{\mathrm p}$ джерело пише одиницю «с$^{-1}$», але за розмірністю \eqref{eq:07-PD} має бути «с». \textbf{Висновок авторів:} за будь-якого розміщення петель деформація зменшує вимірюваний сигнал і послаблює стабілізацію~\cite[с.~28--29]{Ward1992}. Застереження: рівноваги штучно дестабілізовано.
\end{corpusexample}

\begin{corpusexample}[вертикальна мода в проєкті JET, 1975]
Автори проєкту відштовхуються від відомого циліндричного результату: витягнутий шнур зі сталою густиною струму без стінки нестійкий щодо $n=0$, близька провідна стінка його рятує, а на межі стійкості плазма зсувається як ціле вздовж витягнутості~\cite[с.~125--126]{JET1975}. JET від цієї ідеалізації далекий: $R_0/a$ мале, замість суцільної оболонки є окремі котушки PF, профіль струму наперед невідомий. Попереднє чисельне дослідження для плаского струму знайшло нестійкість \emph{навіть} з ідеальною стінкою на зовнішньому контурі розрахункової області; окремі котушки, за оцінкою авторів, лише погіршили б картину~\cite[с.~126]{JET1975}. Інша оцінка, де плазму зсували жорстко, давала стійкість, але вона вводила в оману: знайдена чисельно власна функція на жорсткий зсув зовсім не схожа. Для профілю з нульовою густиною струму на межі нестійкості не знайшли; на випадок, якщо мода все ж з'явиться, проєкт покладався на зворотний зв'язок~\cite[с.~127]{JET1975}. Урок той самий, що й у рамці вище: пасивна стінка вертикальну моду не рятує, а модель «жорсткого шнура» може хибно обіцяти стійкість. Наша 3D-реконструкція JET використовує рівновагу Соловйова з майже пласким струмом (розд.~\ref{ch:jet}); як статична модель геометрії вона придатна, але як модель робочого розряду без керування положенням --- ні.
\end{corpusexample}

\textbf{LQG-керування в RFX-mod}~\cite{Merlo2011}. Інкременти горизонтальної та вертикальної мод лінеаризованої моделі наведено в п.~\ref{sec:04-shape}. Спершу центроїд струму стабілізують два незалежні пропорційні SISO-контури на сідлових котушках, $K_R=\SI{2.0e4}{В/м}$ і $K_Z=\SI{2.6e4}{В/м}$. Домінантна стала часу замкненої системи $\tau_{\mathrm{SISO}}=\SI{209}{мс}$~\cite[с.~91]{Merlo2011}. Потім модель порядку $n=250$ зводять до $n_R=26$ за нормою Ганкеля~\cite[с.~97, 100]{Merlo2011}. Фільтр Калмана має $\tauKF=\SI{3.80}{мс}$~\cite[с.~103]{Merlo2011}, LQR~--- $\tauLQR=\SI{37.3}{мс}$ (у джерелі $\tau_S$, $\tau_F$, $\tau_R$); швидше за $\approx\SI{25}{мс}$ полюс не поставити, бо з'являються коливання~\cite[с.~107]{Merlo2011}. LQG відновлює зазори після збурень $\li$, $\betap$ і тримає відстань до стінки $\gtrsim\SI{3}{см}$~\cite[с.~113]{Merlo2011}. Інтегральної дії немає (після незворотних збурень лишається статична похибка), і все перевірено лише в симуляції.

\textbf{Модель Дуби: критичний розбір.} У~\cite{Duba2018} кожен із трьох каналів «потік (Вб) $\to$ зміщення ($10^{-6}$~м)» описано ланкою $W(p)=1/(p^2+k)$ (змінна Лапласа $p$; у джерелі $W(s)$) з $k=\num{8.5e7}$, \num{9.0e7} і \SI{9.1e7}{с^{-2}} та перехресними підсилювачами порядку $10^5$--$10^7$ (рис.~2 на с.~3). Система без регулятора дає незгасні коливання (рис.~4--6 на с.~4), а з регулятором «P» перехідні процеси загасають за $\lesssim\SI{6}{мс}$ (рис.~7 на с.~4). Ні структуру регулятора, ні походження моделі та коефіцієнтів не розкрито. Головна фізична вада полягає в тому, що полюси $\pm\mathrm{i}\sqrt k$ (частота $\sqrt k/2\pi\approx\SI{1.5}{кГц}$) описують \emph{нейтрально стійкий} осцилятор. Вертикальна нестійкість має дійсний додатний полюс: у жорсткій моделі передавальна функція має вигляд $1/(p^2-\gamma^2)$, а з урахуванням стінки це полюс $p=+\gamma$. Тож модель придатна лише як вправа з простору станів, а не як модель токамака (порівняйте із задачею~2).

\section{Пилкоподібні коливання і поверхня \texorpdfstring{$q=1$}{q=1}}\label{sec:07-saw}

\begin{outofcorpus}[кінк-моди, критерій Крускала--Шафранова, перез'єднання Кадомцева]
Зовнішня кінк-мода $m/n$ нестійка, коли резонансна поверхня $q=m/n$ лежить трохи за межею плазми. Для $m=n=1$ звідси \emph{критерій Крускала--Шафранова} $q(a)>1$. На практиці через моди $m=2,3$ потрібно $\qnf\gtrsim2$--3 (розд.~\ref{ch:trap}). \emph{Внутрішня} мода $m=1$ існує лише тоді, коли в плазмі є поверхня $q=1$, тобто $q(0)<1$. У циліндрі в головному порядку вона маргінальна; у торі знак $\delta W$ і інкремент визначає $\betap$ всередині $q=1$ (Бюссак). \textbf{Цикл пилки:} омічний нагрів загострює профіль струму, $q(0)$ опускається нижче 1, розвивається мода $1/1$, і настає швидкий обвал центральної температури. \textbf{Модель Кадомцева:} гелікальний потік $\psi_*=\psi-\psit$ (на радіан; $\dd\psi_*/\dd\psit=1/q-1$) має екстремум на $q=1$. Поверхні з однаковим $\psi_*$ по обидва боки від $r_1$ перез'єднуються, у зоні $r<r_{\mathrm{mix}}$ (де $\psi_*(r_{\mathrm{mix}})=\psi_*(0)$) профілі вирівнюються, і після обвалу $q\ge1$. Для параболічного струму $r_{\mathrm{mix}}\approx\sqrt2\,r_1$. Час перез'єднання $\tau_K\sim\sqrt{\tau_A\tau_R}$, де $\tau_A=a/\vA$, $\tau_R=\mu_0a^2/\eta$ (резистивний час на всьому радіусі $a$).
\end{outofcorpus}

Велика область з $q<1$ загрожує утриманню α-частинок. У~\cite[с.~4]{Mazzucato2000} нагадано застереження прихильників Ignitor: за $\qnf=3$, як в ITER, такою областю керуватимуть внутрішні моди $m=1$. $\qnf=3{,}6$ можна отримати, збільшивши $a$ на 5\,\% і $\kappa$ до 1,8. Fenix містить редуковані моделі пилки й NTM~\cite[с.~4]{Fable2022}. У базовому сценарії EU-DEMO обвали пилки видно на $T_e(0)$ при $t\approx70$ і \SI{140}{с}~\cite[с.~10]{Fable2022}.

\begin{corpusexample}[придушення пилки ЕЦ-нагрівом на WT-3~\cite{Hanada1990}]
WT-3: $R=\SI{65}{см}$, $a=\SI{20}{см}$, $\Btor\le\SI{1.75}{Тл}$, $\Ip\le\SI{150}{кА}$; гіротрон \SI{56}{ГГц}, $P_{\mathrm{ECH}}\le\SI{200}{кВт}$, X-мода, друга гармоніка. Пилку спостерігали при $q_L=2{,}2$--6,0. Резонанс пересували зміною $\Btor$~\cite[с.~3]{Hanada1990}. \textbf{Локалізація вирішальна.} Нагрів на $q=1$ з боку сильного поля (HFS) повністю пригнічує пилку: період досягає \SI{30}{мс}, тобто тривалості ECH (підпис до рис.~2, с.~9). З боку слабкого поля (LFS) період зростає до $\tausaw=\SI{3}{мс}$, що більше за $\tauE=\SI{1}{мс}$, але повного придушення немає. Нагрів на осі посилює обвали~\cite[с.~3--4]{Hanada1990}. \textbf{Пороги.} При $q_L=3{,}4$, $n_e=\SI{5e18}{м^{-3}}$ (переведено з СГС) і нагріві з HFS повна стабілізація настає при $P_{\mathrm{ECH}}/P_{\mathrm{OH}}\approx3{,}5$. При $q_L\approx4{,}7$ досить $\approx1$, і тоді стабілізує навіть LFS, але з порогом у 1,7 раза вищим~\cite[с.~4--5]{Hanada1990}. Стабілізація можлива лише при $n_e\lesssim(5\text{--}6)\cdot10^{18}$~м$^{-3}$, коли в спектрі жорсткого рентгену з'являються гарячі електрони з $T_{eh}\approx\SI{60}{кеВ}$~\cite[с.~5]{Hanada1990}.
\textbf{Пояснення авторів.} ECH підвищує $T_{e\perp}$ ($\delta T_{e\perp}/T_{e\perp}\approx0{,}3$) і знижує опір. Струм біля $q=1$ зростає, і шир $s=(2/q)(\dd q/\dd\psi)\,V/(\dd V/\dd\psi)$ ($V$~--- об'єм) зменшується~\cite[с.~5]{Hanada1990}. PEST-розрахунок ($q_{\mathrm{axis}}=0{,}83$, $q_L=3{,}37$, $q=1$ на $0{,}35a$) показує, що зі зменшенням ширу рівновага переходить від резистивно нестійкої моди $1/1$ до стійкої~\cite[с.~5--6]{Hanada1990}. HFS/LFS-асиметрію пояснюють захопленням: при нагріві з LFS швидкі електрони потрапляють на бананові орбіти й не дають внеску в струм. Додатково стабілізують швидкі електрони, що прецесують швидше за діамагнітний дрейф, $\omega_{de}>\omega^*_e/2$, де $\omega_{de}=W_\perp/(e\,rR\Btor)$ і $\omega^*_e=(k_\theta/e\Btor)\grad p/n_0$. Звідси $W_\perp>R\grad p/n_0\sim T_eR/a\approx\SI{1.5}{кеВ}$, що набагато менше за 60~кеВ~\cite[с.~6--7]{Hanada1990}. \textbf{Застереження:} $q(0)$ не виміряно; яку саме $P_{\mathrm{OH}}$ підставлено в поріг, не вказано (задача~3).
\end{corpusexample}

\section{Межі за \texorpdfstring{$\beta$}{β} і густиною}\label{sec:07-limits}

\begin{outofcorpus}[межа Троєна і межа Грінвальда]
Ідеальні кінк- і балонні моди обмежують $\beta$. Числові оптимізації дають межу Троєна $\beta_{\max}[\%]=g_{\mathrm{Tr}}\,\Ip[\text{МА}]/(a[\text{м}]\,B[\text{Тл}])$ з $g_{\mathrm{Tr}}\approx2{,}8$ без стінки, тобто $\betaN\lesssim2{,}8$. Часто використовують і $\betaN\lesssim4\li$. Межа Грінвальда емпірична: за $\bar n\to\nGW$ посилюються крайове випромінювання і MARFE, і розряд закінчується зривом (розд.~\ref{ch:disruptions}). Від $\Btor$ і нагріву ця межа майже не залежить.
\end{outofcorpus}

У корпусі обидві величини визначено явно~\cite[с.~2]{Mazzucato2000}:
\begin{equation}\label{eq:07-nGW}
\nGW[10^{20}\,\text{м}^{-3}]=\frac{\Ip[\text{МА}]}{\pi a^2[\text{м}^2]},\qquad n_G=\bar n/\nGW,\qquad \betaN=\frac{\beta}{\Ip/(aB)}
\end{equation}
(у $\betaN$: $\beta$ у~\%, $\Ip$ у МА, $a$ у м, $B$ у Тл). Там же зазначено, що утримання швидко погіршується при $n_G\to1$. $n_G=0{,}85$ ITER-FEAT автор вважає все ще завеликим, а $\betaN$ ITER, за його словами, лежить на верхньому краю бази даних. У табл.~2~\cite[с.~8]{Mazzucato2000} для ITER-FEAT ($B=\SI{5.3}{Тл}$, $\Ip=\SI{15}{МА}$, $a=\SI{2}{м}$, $\beta=2{,}50\,\%$) $\betaN=1{,}77$ і $n_G=0{,}85$, а при $B=\SI{13}{Тл}$ $\betaN=0{,}99$ і $n_G=0{,}39$. Для IAM (табл.~1, с.~7) $\betaN=2{,}25$. \textbf{Балонна межа} в сферичному торі~--- $\beta\approx12\,\%$ у прогнозі TSC+DCON для NSTX~\cite[с.~5]{Jardin2000} (п.~\ref{sec:04-configs}). Її спричиняють малий шир біля $q=1$ і пікований тиск, тобто обидва «важелі» з \eqref{eq:07-dW}.

\begin{practicum}[наскільки далеко DIII-D від меж?]
З кореня проєкту: \texttt{"fusion equilibrium challenge/starter/.venv/bin/python"}\\\texttt{manual/code/ch07\_limits.py}. Для трьох локальних розрядів DIII-D скрипт бере мітки EFIT ($\betaN$, $\li$, $\qnf$), $a$ з LCFS, $\Ip$ з виправленою часовою віссю (\texttt{data\_fixes.fix\_d3d\_ip\_times}; без неї $\Ip$ читається з шуму до пробою) і густину core-Томсона. Далі він рахує $\nGW$ за \eqref{eq:07-nGW}. \textbf{Очікуваний результат (наш прогін):} \#203702/3: $\Ip\approx\SI{1.36}{МА}$, $a\approx\SI{0.605}{м}$, $\qnf\approx3{,}7$, $\betaN\le0{,}42$, $\li\approx1{,}3$, $\max\betaN/(4\li)=0{,}10$, $\langle n_e\rangle_{\mathrm{TS}}/\nGW\approx0{,}10$ (за максимумом по каналах $\le0{,}33$); \#203704: $\qnf\approx5{,}1$, $\betaN\le0{,}58$, $\langle n_e\rangle_{\mathrm{TS}}/\nGW\le0{,}17$ ($\le0{,}42$). Тут $\langle n_e\rangle_{\mathrm{TS}}$~--- середнє по каналах core-Томсона, а не лінійно усереднена густина. Ці розряди далекі від обох меж. \textbf{Завдання:} (1)~пройдіть потоково (\texttt{example\_usage.py}) щонайменше 500 розрядів \texttt{diii\_d\_train} і побудуйте розподіл $\betaN/\li$ та $\langle n_e\rangle_{\mathrm{TS}}/\nGW$; (2)~знайдіть, чи обриваються розряди поблизу $\betaN\approx4\li$ або $\bar n\approx\nGW$; (3)~оцініть, як заміна лінійно усередненої густини середнім по каналах Томсона зсуває $n_G$.
\end{practicum}

\section{ELM і пілінг-балонні моди}\label{sec:07-elm}

\begin{outofcorpus}[пілінг--балонінг]
Градієнт тиску педесталу живить балонні моди, крайовий (бутстреп-, \eqref{eq:05-jbs}) струм~--- кінк-подібні «пілінг»-моди. Їхня спільна межа в площині ($J_{\mathrm{edge}}$, $\grad p$) задає ELM типу~I, що частково руйнує педестал. У редукованих моделях її заміняє обмеження тиску педесталу ($P_{\mathrm{ped,crit}}$ у Fenix, п.~\ref{sec:05-scal}).
\end{outofcorpus}

\begin{corpusexample}[ELM типу~I в ASDEX Upgrade: JOREK~\cite{vanVugt2019}]
Моделюють конвективний ELM \#33616 з гармоніками $n=1\ldots8$~\cite[с.~4]{vanVugt2019}. Спершу нестійка пілінг-балонна мода $n=6$, далі вона зв'язується з $n=5$ та іншими. Домінують $n=3$--5, а $n>6$ залишаються слабкими. ELM забирає $\approx7\,\%$ частинок і $\approx2{,}5\,\%$ енергії за $\approx\SI{2}{мс}$~\cite[с.~6]{vanVugt2019}. Як ці вихори переносять вольфрам, розглянуто в п.~\ref{sec:05-imp} (там же обмеження моделі).
\end{corpusexample}

\section{Тиринг-моди, острови й альфвенівські моди}\label{sec:07-tearing}

\begin{outofcorpus}[$\Delta'$, рівняння Резерфорда, NTM]
Скінченний опір дозволяє розрив і перез'єднання силових ліній на раціональній поверхні $r_s$ ($q=m/n$). Зовнішню ідеальну задачу характеризує стрибок логарифмічної похідної збуреного потоку $\Delta'=[\delta\psi'/\delta\psi]_{r_s-0}^{r_s+0}$, і тиринг-мода нестійка при $\Delta'>0$. У нелінійній фазі острів радіальної ширини $w_r$ росте повільно, на резистивному масштабі $\tau_R=\mu_0r_s^2/\eta$, визначеному на $r_s$ (на відміну від $\mu_0a^2/\eta$ у п.~\ref{sec:07-saw}; рівняння Резерфорда). Для \emph{неокласичної} тиринг-моди (NTM) до нього додається внесок втраченого всередині острова бутстреп-струму:
\begin{equation}\label{eq:07-MRE}
\frac{\tau_R}{r_s}\frac{\dd w_r}{\dd t}\simeq r_s\Delta'(w_r)+C_{\mathrm{bs}}\sqrt\varepsilon\,\betap\frac{L_q}{L_p}\frac{r_sw_r}{w_r^2+w_{\mathrm d}^2},
\end{equation}
де $C_{\mathrm{bs}}=\order{1}$, $L_q$, $L_p$~--- масштаби $q$ і $p$, а $w_{\mathrm d}$~--- ширина, нижче якої поперечна теплопровідність відновлює градієнт. NTM метастабільна: при $\Delta'<0$ вона росте лише із «затравкового» острова ширшого за $\sim w_{\mathrm d}$ (наприклад, після обвалу пилки); її придушують локальною ЕЦ-генерацією струму.
\end{outofcorpus}

Геометрію острова описано в п.~\ref{sec:02-ham}: ширина \eqref{eq:02-w} задана в полоїдальному потоці, у метрах $w_r\approx w_{mn}/(R\Bpol)$. Одна мода дає острів, але не хаос (E20~\cite{ProjRegisters}); стохастичний шар потребує перекриття островів різних мод (обернена задача відновлення спектра мод~--- колишня гіпотеза C5, спростована для базового розв'язувача як R13, розд.~\ref{ch:particles}). Fenix містить редуковану модель NTM (її вигляд у~\cite{Fable2022} не наведено)~\cite[с.~4]{Fable2022}.

\textbf{Альфвенівські моди й острови у фазовому просторі}~\cite{Tyshchenko2021}. У NSTX глобальні альфвенівські моди (GAE) з $\tilde B_r/B\sim10^{-3}$ і $f\approx\SI{0.97}{МГц}$ утворюють резонансні острови у фазовому просторі швидких іонів (аналітика й код розходяться на 20--30\,\%)~\cite[с.~81]{Tyshchenko2021}. Близько десяти таких мод досить, щоб зробити рух стохастичним і відібрати приблизно половину енергії швидких іонів~\cite[с.~119]{Tyshchenko2021}. «Цеберне» перенесення (bucket transport) при гальмуванні іонів біля магнітних островів лише перемішує іони, адвекції воно не створює. Потоки від нього значні за малого ширу й зосереджені між островами~\cite[с.~5, 118]{Tyshchenko2021}. Каналювання енергії~--- розд.~\ref{ch:heating}.

\subsection*{Контрольні запитання}
\begin{enumerate}
\item Які доданки \eqref{eq:07-dW} дестабілізують? Чому балонна мода має $\kpar\approx0$?
\item Що дає \eqref{eq:07-wall} для ідеальної стінки і без стінки? Чому пасивна стінка не стабілізує вертикальну моду остаточно?
\item Чому посилення $g_{\mathrm p}$ з внутрішніми петлями PBX-M не стабілізує плазму, і чому тут не допомагає диференційна складова~\cite{Ward1992}?
\item Чим відрізняються умови $q(a)>1$ і $q(0)\ge1$? Чому ECH з LFS діє слабше, ніж з HFS~\cite{Hanada1990}?
\item Чому NTM метастабільна? Як \eqref{eq:02-w} пов'язана з $w_r$ у \eqref{eq:07-MRE}?
\end{enumerate}

\subsection*{Задачі}
\begin{enumerate}
\item \emph{Межі ITER-FEAT і 13-Тл версії.} За \eqref{eq:07-nGW} і табл.~2~\cite{Mazzucato2000} перевірте $\betaN$ для ITER-FEAT і для $B=\SI{13}{Тл}$ ($\beta=1{,}40\,\%$, $\Ip=\SI{9.7}{МА}$, $a=\SI{0.53}{м}$). Знайдіть $\nGW$ і $\bar n$ в обох випадках, а також запас до межі Троєна $\betaN/2{,}8$. Чи означає менше $n_G$ меншу абсолютну густину?
\item \emph{Нестійкий полюс проти осцилятора.} (а)~За $\tau_{\mathrm{vert}}=\SI{60.72}{мс}$ (п.~\ref{sec:04-shape}) знайдіть $\gamma$ і порівняйте з мінімальним досяжним $\gamma\approx\SI{4}{с^{-1}}$ далеких петель PBX-M. У скільки разів $\tauKF$ і $\tauLQR$~\cite{Merlo2011} коротші за $\tau_{\mathrm{vert}}$? (б)~Для ланки $1/(p^2+k)$~\cite{Duba2018} з $k=\SI{8.5e7}{с^{-2}}$ знайдіть частоту і період і перевірте їх за рис.~4 джерела (горизонт \SI{10}{мс}). (в)~Запишіть передавальну функцію жорсткого шнура з $n_{\mathrm d}<0$ без стінки. Де лежать її полюси?
\item \emph{Поріг ECH на WT-3.} Для умов рис.~3~\cite{Hanada1990} ($\Ip=\SI{80}{кА}$, $\Btor=\SI{0.91}{Тл}$) перевірте $q_L=3{,}4$ за \eqref{eq:01-qcyl}. Оцініть $P_{\mathrm{OH}}=\Vloop\Ip$ при $\Vloop=1{,}5$ і \SI{0.8}{В} (без і з ECH; дані для LFS, $q_L=4{,}3$; для умов рис.~3~--- лише оцінка): за якого $P_{\mathrm{OH}}$ поріг $3{,}5$ сумісний з $P_{\mathrm{ECH}}\le\SI{213}{кВт}$? Порівняйте $T_eR/a$ з $T_{eh}$; оцініть $r_{\mathrm{mix}}$ при $r_1=0{,}35a$.
\end{enumerate}
